\documentclass[12pt,a4paper]{article}
\usepackage[utf8]{inputenc}
\usepackage[spanish]{babel}
\usepackage{geometry}
\geometry{margin=2.5cm}
\usepackage{tabularx}
\usepackage{booktabs}
\usepackage{enumitem}
\usepackage{titlesec}
\usepackage{fancyhdr}

% Configuración de encabezado y pie de página
\pagestyle{fancy}
\fancyhf{}
\setlength{\headheight}{15pt}
\lhead{Sistema Sync-Stock: Informe Final - 3ra Entrega}
\rhead{Junio 2026}
\rfoot{\thepage}

\begin{document}

% --- PORTADA ---
\begin{titlepage}
    \centering
    {\bfseries\Large Universidad Diego Portales \par}
    {\large Facultad de Ingeniería y Ciencias \par}
    {\large Curso Arquitectura de Software \par}
    \vspace{3cm}
    {\bfseries\huge Sistema de Gestión de Inventario Omnicanal \par}
    {\huge (Sync-Stock) \par}
    \vspace{1.5cm}
    {\Large Informe Final - 3ra Entrega \par}
    \vspace{2cm}
    \begin{flushleft}
        {\bfseries Organización:} Tienda de Pesca "PescaTotal" \\
        {\bfseries Área:} Departamento de Ventas e Inventario \\
        {\bfseries Responsables:} Fernando Llancao, Carlos Jara, Bastian Peña, Nicolas Canelo \\
        {\bfseries Fecha:} 7 de Junio de 2026 \\
        {\bfseries Ubicación:} Región Metropolitana, Chile
    \end{flushleft}
    \vfill
    {\large Semestre 1/2026}
\end{titlepage}

\newpage
\tableofcontents
\newpage

% --- 1. EQUIPO DE TRABAJO ---
\section{Equipo de Trabajo}
Para el desarrollo del sistema Sync-Stock, se ha conformado un equipo de trabajo liderado por el responsable del área de desarrollo:

\begin{table}[h]
\centering
\begin{tabularx}{\textwidth}{|l|l|X|}
\hline
\textbf{Integrante} & \textbf{Cargo} & \textbf{Responsabilidades} \\ \hline
Fernando Llancao & CEO / Líder de Proyecto & 
\begin{itemize}[nosep, leftmargin=*]
    \item Liderazgo estratégico y toma de decisiones ejecutivas.
    \item Coordinación directa con los dueños de PescaTotal.
    \item Supervisión general del análisis y diseño del sistema.
\end{itemize} \\ \hline
Nicolas Canelo & Arquitecto de Software & 
\begin{itemize}[nosep, leftmargin=*]
    \item Diseño de la estructura técnica y componentes.
    \item Selección del stack tecnológico adecuado para tiempo real.
\end{itemize} \\ \hline
Carlos Jara & Analista de Sistemas & 
\begin{itemize}[nosep, leftmargin=*]
    \item Levantamiento y especificación de requerimientos.
    \item Definición del modelo de datos y persistencia.
\end{itemize} \\ \hline
Bastian Peña & Desarrollador Full-stack & 
\begin{itemize}[nosep, leftmargin=*]
    \item Implementación de la lógica de negocio y API.
    \item Desarrollo de la interfaz de usuario amigable.
\end{itemize} \\ \hline
\end{tabularx}
\end{table}

% --- 2. DESCRIPCIÓN DEL SISTEMA Y CONTEXTO ---
\section{Descripción del Sistema y Contexto Organizacional}

\subsection{La Organización: PescaTotal}
PescaTotal es una tienda retail especializada en artículos de pesca recreativa y equipamiento outdoor. A pesar de ser una organización de escala pequeña, posee un flujo de ventas constante y una base de clientes fidelizada en la Región Metropolitana.

\subsection{Área de Trabajo: Departamento de Ventas e Inventario}
Esta unidad operativa gestiona el stock físico de la tienda y, recientemente, la administración del nuevo canal e-commerce. Sus funciones incluyen la recepción de suministros, atención al público y preparación de pedidos online.

\subsection{Problemática Actual}
La tienda enfrenta dificultades operacionales críticas derivadas de la falta de integración digital:
\begin{itemize}
    \item \textbf{Incumplimiento de pedidos y fricción con clientes por falta de gestión de reservas:} Al carecer de un mecanismo que bloquee o separe físicamente los productos apenas se venden por el canal e-commerce, los artículos quedan expuestos en las vitrinas y son comprados por clientes presenciales. Esto genera un problema grave para el negocio: la necesidad de cancelar compras online ya pagadas, procesar devoluciones de dinero, sobrecargar al personal de atención al cliente y generar un impacto negativo en la reputación y fidelidad de los compradores de PescaTotal.
\end{itemize}

\subsection{Sistema Propuesto: Sync-Stock}
Sync-Stock es una plataforma segura de gestión de inventario omnicanal. Su objetivo principal es actuar como un núcleo de datos centralizado que sincronice el stock en tiempo real entre la tienda física y el e-commerce.

Para garantizar la integridad de las operaciones, el sistema contará con acceso restringido mediante roles, entregando a cada empleado (Vendedor, Encargado E-commerce o Administrador) solo las herramientas que necesita. Al registrar una venta en el mesón o recibir un pedido web, el sistema no solo actualiza la disponibilidad global inmediatamente, sino que genera un estado de "Reserva" inmutable para las compras web, bloqueando la posibilidad de sobreventas en la tienda física. Esto asegurará la coherencia total de los datos y protegerá la experiencia de compra de los clientes.

% --- 3. OBJETIVOS Y PERFILES ---
\section{Objetivos y Perfiles de Usuario}

\subsection{Objetivos del Sistema}
\subsubsection{Objetivo General}
Implementar un sistema integral que garantice la coherencia total del inventario entre la tienda física y el e-commerce de PescaTotal, optimizando la eficiencia operativa y eliminando errores de sobreventa.

\subsubsection{Objetivos Específicos}
\begin{enumerate}
    \item Sincronizar de forma automática y bidireccional el stock tras cada transacción.
    \item Implementar un módulo de reserva inmediata para asegurar productos vendidos vía web.
    \item Proporcionar visibilidad de stock real al personal de ventas mediante una interfaz simple.
    \item "Reducir en un 85\% las horas-hombre semanales dedicadas a la conciliación y revisión manual de inventarios, alcanzando esta métrica al finalizar el segundo mes de operación en producción del sistema Sync-Stock."
\end{enumerate}

\subsection{Perfiles de Usuario}
\begin{itemize}
    \item \textbf{Usuario Primario (Vendedor):} Utiliza el sistema en el mesón de la tienda para consultar disponibilidad y cerrar ventas presenciales.
    \begin{itemize}
        \item Frecuencia de uso: Permanente durante la jornada.
        \item Nivel tecnico: Intermedio.
    \end{itemize}
    \item \textbf{Usuario Secundario (Encargado E-commerce):} Gestiona los pedidos web, visualiza las reservas y confirma la salida física del producto.
    \begin{itemize}
        \item Frecuencia de uso: Permanente durante la jornada.
        \item Nivel tecnico: Intermedio.
    \end{itemize}
    \item \textbf{Usuario de Apoyo (Administrador):} Gestiona el catálogo, permisos de usuarios y configuraciones de alertas de stock.
    \begin{itemize}
        \item Frecuencia de uso: Variable, segun necesidades de la tienda.
        \item Nivel tecnico: Avanzado.
    \end{itemize}
\end{itemize}

% --- 4. REQUERIMIENTOS FUNCIONALES ---
\section{Requerimientos Funcionales}

\begin{table}
\centering
\begin{tabularx}{\textwidth}{|l|l|X|}
\hline
\textbf{Código} & \textbf{Nombre} & \textbf{Descripción} \\ \hline
RF-001 & Sincronización Real-Time & El sistema debe descontar el stock del catálogo online inmediatamente tras una venta física confirmada. \\ \hline
RF-002 & Reserva de Stock Online & Al detectarse una orden de compra web, el sistema debe mover el producto al estado "Reservado", restándolo del stock físico vendible. \\ \hline
RF-003 & Gestión de Productos & Permitir el registro, edición y categorización de artículos con códigos de barra y ubicación física. \\ \hline
RF-004 & Control de Despachos & Funcionalidad para confirmar la salida de productos reservados, actualizando el estado a "Vendido" definitivamente. \\ \hline
RF-005 & Alertas de Quiebre & Notificaciones automáticas cuando el stock disponible llegue a un umbral mínimo configurado. \\ \hline
RF-006 & Log de Auditoría & Registro detallado de quién, cuándo y por qué canal se realizó cada movimiento de inventario. \\ \hline
RF-007 & Autenticación de Usuarios (Login) & El sistema debe requerir que todo usuario inicie sesión mediante credenciales (correo/usuario y contraseña) antes de acceder a cualquier funcionalidad. \\ \hline
RF-008 & Control de Acceso Basado en Roles (RBAC) & El sistema debe restringir las funcionalidades e interfaces de acuerdo con el rol del usuario (Administrador, Vendedor, Encargado E-commerce). Por ejemplo, solo el Administrador podrá gestionar permisos o el catálogo. \\ \hline
RF-009 & Cierre de Sesión y Tiempos de Inactividad & El sistema debe permitir el cierre de sesión manual y forzar un cierre automático tras 30 minutos de inactividad por motivos de seguridad en los equipos compartidos de la tienda. \\ \hline
\end{tabularx}
\end{table}

\newpage

\section{Modelo de Datos y Mecanismo de Persistencia}
Para dar soporte a las funcionalidades del sistema Sync-Stock en un entorno retail omnicanal real, se ha diseñado una estrategia de persistencia transaccional de alta consistencia y un modelo de datos relacional detallado.

\subsection{Mecanismo de Persistencia y Decisiones de Diseño}
Se ha optado por implementar una \textbf{Base de Datos Relacional (RDBMS)} utilizando \textbf{PostgreSQL}. Esta decisión responde a la necesidad de resolver problemas complejos de concurrencia y consistencia que un sistema omnicanal enfrenta en tiempo real:

\begin{enumerate}
    \item \textbf{Consistencia Transaccional Estricta (ACID):} Para evitar la sobreventa, los cambios en el inventario deben procesarse con aislamiento transaccional estricto. PostgreSQL permite aislar operaciones mediante el nivel de aislamiento de transacciones \textit{Read Committed} combinado con bloqueos pesimistas explícitos.
    
    \item \textbf{Estrategia de Concurrencia y Bloqueos Pesimistas (Pessimistic Locking):} Al procesar una compra en el Punto de Venta (POS) o iniciar un proceso de reserva web, el sistema ejecuta una transacción que bloquea las filas correspondientes de la tabla de inventario mediante la sentencia \texttt{SELECT ... FOR UPDATE}.
    \begin{itemize}
        \item Esto garantiza que ningún otro proceso pueda modificar la cantidad del producto hasta que la transacción actual finalice (haga \textit{COMMIT} o \textit{ROLLBACK}).
        \item \textbf{Prevención de Deadlocks:} Para evitar bloqueos mutuos (\textit{deadlocks}) bajo cargas de alta concurrencia, el sistema ordena siempre los ítems de la consulta de forma determinista (ordenados por el identificador del producto) antes de solicitar el bloqueo.
    \end{itemize}
    
    \item \textbf{Política de Expiración de Reservas y Liberación de Stock (Time-To-Live):} El canal e-commerce genera reservas temporales para proteger el stock de un cliente mientras completa el pago. Sin embargo, para evitar que el stock quede bloqueado indefinidamente si el cliente abandona el carrito, se implementa una regla de expiración de 15 minutos.
    \begin{itemize}
        \item Cada reserva se registra en la tabla \texttt{Reserva} con un campo \texttt{FechaExpiracion}.
        \item Un proceso daemon (o \textit{cron job}) en el bus de servicios ejecuta periódicamente (cada 1 minuto) una transacción que selecciona las reservas expiradas en estado \texttt{Pendiente}, revierte el stock reservado al stock disponible (\texttt{StockDisponible = StockDisponible + Cantidad}) y marca la reserva como \texttt{Expirada}.
    \end{itemize}
    
    \item \textbf{Pooling de Conexiones y Alta Disponibilidad:} Para asegurar tiempos de respuesta menores a 200 ms y soportar picos de carga concurrentes entre la tienda y el sitio web, se utiliza un pooler de conexiones (\textbf{PgBouncer}) configurado en modo transacción. Esto permite optimizar el uso de recursos y mantener la latencia en rangos óptimos.
\end{enumerate}

\subsection{Diccionario de Datos (Entidades del Sistema Omnicanal)}
A continuación se detalla el esquema de base de datos extendido que soporta la lógica del inventario multicanal, reservas y auditoría de transacciones:

\begin{table}[h!]
\centering
\small
\begin{tabularx}{\textwidth}{@{}lllX@{}}
\toprule
\textbf{Entidad} & \textbf{Atributo} & \textbf{Tipo de Dato} & \textbf{Descripción} \\ \midrule
\textbf{Producto} & ID\_Producto (PK) & UUID & Identificador universal único del artículo. \\
 & SKU & VARCHAR(50) & Código SKU único (Stock Keeping Unit). \\
 & Nombre & VARCHAR(150) & Nombre o descripción del artículo. \\
 & Precio & DECIMAL(12,2) & Precio de venta unitario. \\ \addlinespace
\textbf{Ubicacion} & ID\_Ubicacion (PK) & UUID & Identificador único de la ubicación de stock. \\
 & Nombre & VARCHAR(100) & Nombre (ej. 'Bodega Central', 'POS Tienda'). \\
 & Tipo & ENUM & Tipo de instalación: 'Tienda' o 'Bodega'. \\ \addlinespace
\textbf{Inventario} & ID\_Inventario (PK) & UUID & Identificador del registro de inventario. \\
 & Producto\_ID (FK) & UUID & Referencia a la entidad \textbf{Producto}. \\
 & Ubicacion\_ID (FK) & UUID & Referencia a la entidad \textbf{Ubicacion}. \\
 & StockDisponible & INT & Unidades físicas disponibles para venta inmediata. \\
 & StockReservado & INT & Unidades retenidas temporalmente por e-commerce. \\
 & StockTransito & INT & Unidades en proceso de traslado entre bodegas. \\ \bottomrule
\end{tabularx}
\caption{Diccionario de datos: Catálogo, Ubicaciones e Inventario.}
\end{table}

\begin{table}[h!]
\centering
\small
\begin{tabularx}{\textwidth}{@{}lllX@{}}
\toprule
\textbf{Entidad} & \textbf{Atributo} & \textbf{Tipo de Dato} & \textbf{Descripción} \\ \midrule
\textbf{Reserva} & ID\_Reserva (PK) & UUID & Identificador único de la reserva temporal. \\
 & Producto\_ID (FK) & UUID & Referencia a la entidad \textbf{Producto}. \\
 & Ubicacion\_ID (FK) & UUID & Referencia a la entidad \textbf{Ubicacion}. \\
 & Cantidad & INT & Cantidad de unidades bloqueadas. \\
 & Estado & ENUM & Estado de la reserva: 'Pendiente', 'Confirmada', 'Expirada'. \\
 & FechaCreacion & TIMESTAMP & Marca de tiempo de la creación de la reserva. \\
 & FechaExpiracion & TIMESTAMP & Límite temporal para consolidar la transacción. \\ \addlinespace
\textbf{Orden} & ID\_Orden (PK) & UUID & Identificador único de la orden de venta. \\
 & Canal & ENUM & Origen del pedido: 'POS\_Fisico' o 'E\_Commerce'. \\
 & Estado & ENUM & Estado del pedido: 'Creada', 'Pagada', 'Despachada', 'Cancelada'. \\
 & FechaCreacion & TIMESTAMP & Marca de tiempo de registro de la orden. \\
 & Total & DECIMAL(12,2) & Monto total de la transacción comercial. \\ \addlinespace
\textbf{LineaOrden} & ID\_Linea (PK) & UUID & Identificador de la línea de detalle de la orden. \\
 & Orden\_ID (FK) & UUID & Referencia a la entidad \textbf{Orden}. \\
 & Producto\_ID (FK) & UUID & Referencia a la entidad \textbf{Producto}. \\
 & Cantidad & INT & Cantidad de unidades compradas. \\
 & PrecioUnitario & DECIMAL(12,2) & Precio facturado para este producto. \\ \bottomrule
\end{tabularx}
\caption{Diccionario de datos: Reservas temporales y Órdenes de venta.}
\end{table}

\begin{table}[h!]
\centering
\small
\begin{tabularx}{\textwidth}{@{}lllX@{}}
\toprule
\textbf{Entidad} & \textbf{Atributo} & \textbf{Tipo de Dato} & \textbf{Descripción} \\ \midrule
\textbf{HistorialInv} & ID\_Historial (PK) & UUID & Identificador único para auditoría de movimientos. \\
 & Producto\_ID (FK) & UUID & Referencia al producto afectado. \\
 & Ubicacion\_ID (FK) & UUID & Ubicación física del movimiento. \\
 & Cantidad & INT & Unidades modificadas (valor positivo o negativo). \\
 & TipoMovimiento & ENUM & Tipo de evento: 'Venta', 'Reserva\_Creada', 'Reserva\_Expirada', 'Recepcion', 'Ajuste'. \\
 & FechaHora & TIMESTAMP & Fecha y hora exacta de la transacción. \\
 & Usuario\_ID (FK) & UUID & Identificador del operario responsable. \\ \bottomrule
\end{tabularx}
\caption{Diccionario de datos: Historial de inventario para log de auditoría.}
\end{table}
\clearpage

\section{Arquitectura SOA y Mediación a través de un Bus de Servicios (ESB)}
Para resolver la integración y sincronización en tiempo real de múltiples canales, el sistema se ha estructurado bajo el patrón \textbf{SOA (Service-Oriented Architecture)} incorporando un \textbf{Enterprise Service Bus (ESB)} central como columna vertebral de comunicación. El ESB centraliza las interacciones y proporciona mecanismos de mediación clásicos que desacoplan físicamente y lógicamente a los clientes y servicios, operando de manera exclusiva en la capa de transporte mediante conexiones TCP persistentes.

\subsection{El Rol del Enterprise Service Bus (ESB) y Límite Arquitectónico}
A fin de asegurar el cumplimiento de los estándares de una arquitectura SOA tradicional y evitar el uso de protocolos acoplados a la web dentro del bus de servicios, se ha establecido un límite arquitectónico estricto. **No existe presencia de HTTP ni de APIs REST dentro del núcleo de la arquitectura SOA.** La comunicación entre el ESB y los microservicios proveedores, así como con los consumidores nativos, se realiza de forma exclusiva mediante sockets TCP utilizando un protocolo de mensajería con entramado de longitud fija (\textit{Length-Prefixed Framing}).

El ESB proporciona las siguientes capacidades de infraestructura:
\begin{itemize}
    \item \textbf{Garantía de Entrega y Asincronía:} Permite encolar las peticiones en canales de mensajería TCP dedicados. Si un servicio (como el de inventario o despacho) experimenta una indisponibilidad temporal, los mensajes TCP quedan retenidos en el bus en un buffer de conexión hasta que el servicio correspondiente restablezca su socket y procese el mensaje, garantizando que no se pierdan transacciones de negocio.
    \item \textbf{Desacoplamiento Tecnológico y de Protocolos:} Al independizar la capa de servicios del protocolo HTTP, la arquitectura SOA puede integrar de forma nativa microservicios construidos en diversos lenguajes y plataformas (como scripts Python de automatización o servicios de backend en Node.js) que se comunican directamente a través de flujos TCP estructurados (JSON/TCP), eliminando la sobrecarga y la dependencia de servidores web tradicionales.
    \item \textbf{Gateway/Adaptador de Clientes Web Externos (Perímetro):} Debido a las restricciones de seguridad de los navegadores web modernos, las aplicaciones frontend en ejecución en el cliente (como React) no pueden abrir sockets TCP crudos directamente. Para solventar esto, se posiciona un **API Gateway / Adaptador de Clientes (Express.js)** en la frontera exterior del sistema. Este adaptador actúa como un cliente externo a la red SOA: recibe las llamadas HTTP REST del navegador web, las valida y autoriza (RBAC mediante tokens JWT) y las traduce a mensajes en formato JSON/TCP que inyecta en el ESB. De esta forma, el protocolo HTTP queda estrictamente relegado al perímetro de red externa, fuera del núcleo de servicios SOA.
\end{itemize}

\subsection{Mecanismos Clásicos de Mediación SOA Implementados}
El ESB actúa como una capa de mediación activa sobre los flujos TCP de datos, aplicando patrones de integración empresarial (EIP) clásicos para procesar los mensajes en tránsito:

\begin{enumerate}
    \item \textbf{Enrutamiento Basado en Contenido (Content-Based Router):} Cuando el ESB recibe una trama TCP de orden o consulta, inspecciona el campo \texttt{action} o \texttt{canal} del payload. 
    \begin{itemize}
        \item Si el mensaje proviene de un canal digital como \texttt{E\_Commerce}, el ESB enruta dinámicamente el mensaje hacia el servicio de inventario para la validación de reservas y flujo de pago.
        \item Si el mensaje proviene de un canal físico como el punto de venta \texttt{POS\_Fisico}, es enrutado directamente al servicio de inventario para una rebaja directa sin estado intermedio de reserva.
    \end{itemize}
    
    \item \textbf{Transformación y Traducción de Datos (Data Transformer):} Para homogeneizar los formatos de las peticiones que ingresan al bus (como payloads provenientes de terminales POS antiguas o plataformas e-commerce con campos abreviados), el mediador aplica un transformador que normaliza los campos de entrada al esquema estandarizado del bus (\textit{Modelo Canónico}) antes de que el mensaje llegue al servicio de inventario.
    
    \item \textbf{Enriquecimiento de Mensajes (Message Enricher):} Cuando un cliente envía una trama TCP simplificada (con el SKU y cantidad únicamente), el bus intercepta el mensaje y consulta internamente al catálogo para enriquecer el payload agregando el nombre del producto, descripción y precio unitario vigente. Esto asegura que los servicios core reciban mensajes completos y contextualizados para su procesamiento.
    
    \item \textbf{Puente de Protocolos (Protocol Bridge) de Frontera:} A nivel perimetral, la pasarela de servicios traduce las peticiones HTTP REST síncronas del frontend web en tramas estructuradas TCP enviadas al bus. Este componente de frontera permite la interoperabilidad de la interfaz de usuario con la arquitectura SOA interna, traduciendo formatos, cabeceras y parámetros HTTP a sockets de red TCP y viceversa.
\end{enumerate}

\subsection{Componentes de la Arquitectura}
Bajo este esquema de mediación, los componentes del ecosistema SOA se dividen en:

\subsubsection{Componentes Cliente (Consumidores de Servicios)}
\begin{itemize}[leftmargin=*]
    \item \textbf{Clientes Directos TCP (POS de Escritorio / Scripts de Automatización):} Clientes de escritorio y procesos automatizados que se conectan directamente al ESB en el puerto 5000 utilizando sockets TCP crudos, consumiendo servicios de manera nativa sin intermediarios HTTP.
    \item \textbf{API Gateway de Frontera (Adaptador REST-TCP):} Servidor web Express que actúa como puente de frontera. Traduce las llamadas HTTP REST del navegador web (Frontend POS, Dashboard de Administración) a mensajes de red TCP compatibles con el bus SOA.
\end{itemize}

\subsubsection{Componentes de Servicio (Proveedores del Bus)}
\begin{itemize}[leftmargin=*]
    \item \textbf{Auth Service (Servicio 'autho'):} Servicio de autenticación registrado en el bus TCP. Procesa credenciales de usuario y genera firmas criptográficas JWT.
    \item \textbf{Inventory Service (Servicio 'inven'):} Servicio core de inventario. Resuelve consultas y realiza reservas o confirmaciones de stock interactuando de forma atómica con la base de datos PostgreSQL.
    \item \textbf{Sales \& Order Service (Servicio 'sales'):} Gestiona el flujo de órdenes, coordinando con el servicio de inventario a través de mensajes TCP y auditando las transacciones de negocio.
\end{itemize}

\section{Interfaces de los Componentes y Contratos de Mensajería TCP}
Para garantizar la interoperabilidad del sistema Sync-Stock en un entorno omnicanal crítico en tiempo real, se definen contratos de mensajería TCP rigurosos a nivel del bus de servicios (ESB). Cada contrato especifica el identificador del servicio destino, el tipo de acción a ejecutar, la estructura de la carga útil del mensaje (request) y la estructura de la respuesta (response). A continuación, se resume la relación entre mensajes, Requerimientos Funcionales (RF) y Requerimientos No Funcionales (RNF):

\begin{table}[h!]
\centering
\footnotesize
\begin{tabularx}{\textwidth}{|l|l|l|X|X|}
\hline
\textbf{Servicio Destino} & \textbf{Acción (Campo \texttt{action})} & \textbf{RF Satisfecho} & \textbf{RNF Mapeado} & \textbf{Protocolo y Puerto} \\ \hline
\texttt{autho} & \textit{Login (Implícito)} & RF-007, RF-008 & Seguridad (RBAC), JWT & TCP Socket, Puerto 5000 \\ \hline
\texttt{inven} & \texttt{reserve} & RF-002 & Consistencia Estricta, Bloqueo Pesimista & TCP Socket, Puerto 5000 \\ \hline
\texttt{inven} & \texttt{commit} & RF-004, RF-006 & Integridad Transaccional, Auditoría & TCP Socket, Puerto 5000 \\ \hline
\texttt{inven} & \texttt{stock} & RF-001, RF-003 & Baja Latencia ($<200$ ms) & TCP Socket, Puerto 5000 \\ \hline
\end{tabularx}
\caption{Resumen de contratos de mensajería TCP de la arquitectura SOA.}
\end{table}

\subsection{Detalle de Contratos de Mensajería TCP}
A continuación se especifican las estructuras de las tramas de petición y de respuesta enviadas a través del ESB en formato JSON mediante sockets TCP crudos. Cabe destacar que todas las respuestas de los servicios inician con un prefijo de estado de dos caracteres (\texttt{OK} o \texttt{NK}) para indicar el éxito o fracaso a nivel del bus de servicios, seguido de la carga útil correspondiente.

\subsubsection{1. Servicio de Autenticación (Servicio \texttt{autho})}
\textbf{Identificador en el Bus:} \texttt{autho} \\
\textbf{Descripción:} Valida las credenciales del operario y devuelve un token firmado (JWT) con su rol asociado para habilitar el control de acceso en los clientes de frontera.

\paragraph{Estructura del Mensaje de Petición (Request TCP Payload):}
\begin{verbatim}
{
  "username": "fllancao",
  "password": "mi_clave_segura_123"
}
\end{verbatim}

\paragraph{Mensaje de Respuesta Exitoso (TCP Response Payload - Prefijo \texttt{OK}):}
La respuesta es enviada por el socket con el código de estado \texttt{OK} prependido en la trama:
\begin{verbatim}
OK{
  "token": "eyJhbGciOiJIUzI1NiIsInR5...[truncado]",
  "rol": "Administrador",
  "username": "fllancao"
}
\end{verbatim}

\paragraph{Mensaje de Respuesta de Error (TCP Response Payload - Prefijo \texttt{NK}):}
En caso de fallo en las credenciales, el bus retorna la señal de error \texttt{NK} y la causa del rechazo:
\begin{verbatim}
NKCREDENTIALS_INVALID
\end{verbatim}

\newpage
\subsubsection{2. Reserva Temporal de Stock (Servicio \texttt{inven}, Acción \texttt{reserve})}
\textbf{Identificador en el Bus:} \texttt{inven} \\
\textbf{Acción:} \texttt{reserve} \\
\textbf{Descripción:} Adquiere un bloqueo pesimista en base de datos y reserva temporalmente unidades de un producto en una ubicación específica para evitar la sobreventa.

\paragraph{Estructura del Mensaje de Petición (Request TCP Payload):}
\begin{verbatim}
{
  "action": "reserve",
  "sku": "PES-CA1",
  "cantidad": 2,
  "ubicacion": "Tienda Valdivia",
  "usuarioId": "e4b2d288-43ea-46b2-a5e2-88229988ff22"
}
\end{verbatim}

\paragraph{Mensaje de Respuesta Exitoso (TCP Response Payload - Prefijo \texttt{OK}):}
\begin{verbatim}
OK{
  "reservaId": 12,
  "estado": "Pendiente",
  "expiraEn": "2026-06-03T01:25:00.000Z"
}
\end{verbatim}

\paragraph{Mensaje de Respuesta de Error (TCP Response Payload - Prefijo \texttt{NK}):}
\begin{verbatim}
NKSTOCK_INSUFICIENTE
\end{verbatim}

\subsubsection{3. Confirmación de Reserva (Servicio \texttt{inven}, Acción \texttt{commit})}
\textbf{Identificador en el Bus:} \texttt{inven} \\
\textbf{Acción:} \texttt{commit} \\
\textbf{Descripción:} Transforma una reserva temporal en una venta confirmada definitiva, rebajando definitivamente el stock reservado del catálogo.

\paragraph{Estructura del Mensaje de Petición (Request TCP Payload):}
\begin{verbatim}
{
  "action": "commit",
  "reservaId": "b1b88220-432a-4632-a5e2-88229988ff22",
  "usuarioId": "e4b2d288-43ea-46b2-a5e2-88229988ff22"
}
\end{verbatim}

\paragraph{Mensaje de Respuesta Exitoso (TCP Response Payload - Prefijo \texttt{OK}):}
\begin{verbatim}
OK{
  "id": "c8e228c2-482a-466d-88b1-2299aaffd883",
  "total": 119980.00,
  "estado": "Pagada"
}
\end{verbatim}

\paragraph{Mensaje de Respuesta de Error (TCP Response Payload - Prefijo \texttt{NK}):}
\begin{verbatim}
NKRESERVATION_NOT_FOUND
\end{verbatim}

\subsubsection{4. Consulta de Stock en Tiempo Real (Servicio \texttt{inven}, Acción \texttt{stock})}
\textbf{Identificador en el Bus:} \texttt{inven} \\
\textbf{Acción:} \texttt{stock} \\
\textbf{Descripción:} Devuelve la disponibilidad detallada de un SKU en todas las ubicaciones físicas y bodegas registradas.

\paragraph{Estructura del Mensaje de Petición (Request TCP Payload):}
\begin{verbatim}
{
  "action": "stock",
  "sku": "PES-CA1"
}
\end{verbatim}

\paragraph{Mensaje de Respuesta Exitoso (TCP Response Payload - Prefijo \texttt{OK}):}
\begin{verbatim}
OK{
  "sku": "PES-CA1",
  "nombre": "Caña de Pescar Shakespeare Ugly Stik",
  "precio": 59990.00,
  "inventario": [
    {
      "bodega": "Tienda Valdivia",
      "tipo": "Tienda",
      "disponible": 10,
      "reservado": 0,
      "enTransito": 0
    },
    {
      "bodega": "Bodega E-commerce",
      "tipo": "Bodega",
      "disponible": 30,
      "reservado": 0,
      "enTransito": 0
    }
  ]
}
\end{verbatim}


\section{Documentación de Servicios y Procesos de Mensajería TCP}

Para soportar las nuevas funcionalidades e integraciones de la 3ra Entrega, se ha extendido el catálogo de contratos de mensajería TCP del sistema Sync-Stock. Las peticiones externas son recibidas por el API Gateway perimetral de Express (actuando como un Protocol Bridge externo), el cual decodifica y valida las credenciales y el token JWT, evalúa el control de acceso por rol (RBAC) con el middleware \texttt{verificarToken} y \texttt{permitirRoles}, empaqueta el contenido y lo retransmite al Enterprise Service Bus (ESB) por sockets TCP en formato JSON.

A continuación se detallan las operaciones adicionales incorporadas al catálogo de servicios del bus TCP:

\subsection{Registro de Nuevos Productos}
\textbf{Servicio Destino:} \texttt{sales} \\
\textbf{Acción (Campo \texttt{action}):} \texttt{create-product} \\
\textbf{Descripción:} Registra un nuevo producto en el catálogo y de forma automática crea un registro de inventario con stock inicial en todas las ubicaciones y bodegas disponibles bajo una única transacción ACID.

\paragraph{Estructura del Mensaje de Petición (Request TCP Payload):}
\begin{verbatim}
{
  "action": "create-product",
  "sku": "PES-SE-NZ-023",
  "nombre": "Señuelo Rapala Magnum",
  "precio": 12500,
  "initialStock": 20
}
\end{verbatim}

\paragraph{Mensaje de Respuesta Exitoso (TCP Response Payload - Prefijo \texttt{OK}):}
\begin{verbatim}
OK{
  "id": "f5b88220-432a-4632-a5e2-88229988ffaa",
  "sku": "PES-SE-NZ-023",
  "nombre": "Señuelo Rapala Magnum",
  "precio": 12500
}
\end{verbatim}

\paragraph{Mensaje de Respuesta de Error (TCP Response Payload - Prefijo \texttt{NK}):}
\begin{verbatim}
NKPRODUCT_ALREADY_EXISTS
\end{verbatim}

\subsection{Modificación de Producto}
\textbf{Servicio Destino:} \texttt{sales} \\
\textbf{Acción (Campo \texttt{action}):} \texttt{update-product} \\
\textbf{Descripción:} Actualiza los atributos básicos (nombre y precio) de un producto existente identificado por su SKU en el catálogo.

\paragraph{Estructura del Mensaje de Petición (Request TCP Payload):}
\begin{verbatim}
{
  "action": "update-product",
  "sku": "PES-SE-NZ-023",
  "nombre": "Señuelo Rapala Magnum Pro",
  "precio": 13990
}
\end{verbatim}

\paragraph{Mensaje de Respuesta Exitoso (TCP Response Payload - Prefijo \texttt{OK}):}
\begin{verbatim}
OK{
  "sku": "PES-SE-NZ-023",
  "nombre": "Señuelo Rapala Magnum Pro",
  "precio": 13990
}
\end{verbatim}

\paragraph{Mensaje de Respuesta de Error (TCP Response Payload - Prefijo \texttt{NK}):}
\begin{verbatim}
NKPRODUCT_NOT_FOUND
\end{verbatim}

\subsection{Despacho de Órdenes Consolidadas}
\textbf{Servicio Destino:} \texttt{sales} \\
\textbf{Acción (Campo \texttt{action}):} \texttt{dispatch-order} \\
\textbf{Descripción:} Marca una orden pagada y reservada temporalmente como despachada físicamente, actualizando su estado interno y liberando la reserva del flujo activo.

\paragraph{Estructura del Mensaje de Petición (Request TCP Payload):}
\begin{verbatim}
{
  "action": "dispatch-order",
  "orderId": "d3b07384-d113-49c5-a5d6-d68aee460183"
}
\end{verbatim}

\paragraph{Mensaje de Respuesta Exitoso (TCP Response Payload - Prefijo \texttt{OK}):}
\begin{verbatim}
OK{
  "id": "d3b07384-d113-49c5-a5d6-d68aee460183",
  "canal": "E_Commerce",
  "estado": "Despachada",
  "fechaCreacion": "2026-06-03T19:30:00.000Z",
  "total": 12500
}
\end{verbatim}

\paragraph{Mensaje de Respuesta de Error (TCP Response Payload - Prefijo \texttt{NK}):}
\begin{verbatim}
NKORDER_NOT_FOUND
\end{verbatim}

\subsection{Redirección General a través del Bus de Servicios (Acciones de Lectura)}
Para dar cumplimiento estricto con el diseño arquitectónico SOA definido para PescaTotal, se definieron contratos TCP específicos para la consulta de información de solo lectura en el servicio \texttt{sales}. Cuando el API Gateway perimetral recibe las peticiones HTTP GET correspondientes, realiza la delegación enviando tramas TCP con las siguientes acciones al bus de servicios:
\begin{itemize}
    \item \textbf{Consulta de Catálogo (\texttt{action: 'products'}):} Recupera la lista completa de productos junto con el stock disponible y reservado por ubicación.
    \item \textbf{Historial de Órdenes (\texttt{action: 'orders'}):} Recupera las últimas 15 órdenes de venta registradas en el sistema.
    \item \textbf{Listado de Reservas (\texttt{action: 'reservations'}):} Recupera el estado de las últimas 10 reservas de stock realizadas por el e-commerce.
    \item \textbf{Auditoría de Movimientos (\texttt{action: 'history'}):} Recupera los últimos 10 eventos de inventario registrados para control de auditoría.
\end{itemize}
De esta forma, el servidor Express actúa únicamente como un puente de red y traductor perimetral, mientras que el procesamiento y acceso a datos se centraliza en el bus de servicios TCP, manteniendo el desacoplamiento físico y de persistencia de la capa de negocio.

\newpage
\section{Documentación del Sistema y Evidencias de Funcionamiento}

\subsection{Explicación Técnica del Código e Implementación}

El sistema Sync-Stock se implementó utilizando Node.js, TypeScript y Prisma, logrando los objetivos mediante los siguientes componentes core:

\begin{enumerate}
    \item \textbf{Bus de Servicios TCP (ESB Customizado):}
    Implementado en \texttt{soabus.ts}, este módulo levanta un socket TCP en el puerto 5000. Los microservicios de inventario (\texttt{inven}), ventas (\texttt{sales}) y autenticación (\texttt{autho}) se registran abriendo sockets persistentes. La mensajería interna utiliza un protocolo de entramado customizado de longitud fija (Length-Prefixed Framing). Cada mensaje enviado posee un prefijo de 5 caracteres numéricos que especifican el tamaño exacto del contenido en bytes, seguido de 5 caracteres que identifican al servicio de destino. Esto previene problemas clásicos de fragmentación de buffer en streams TCP socket concurrentes.
    
    \item \textbf{Control de Concurrencia y Bloqueos Pesimistas:}
    Para garantizar que no ocurran sobreventas en checkout simultáneos (POS físico y web), el servicio de inventario (\texttt{inventory.ts}) utiliza transacciones pesimistas (\texttt{prisma.\$transaction}). Dentro de la transacción, se ejecuta una consulta raw SQL con \texttt{SELECT * FROM "Inventario" ... FOR UPDATE}. Esto bloquea temporalmente las filas de stock específicas de los productos requeridos, forzando a que cualquier transacción concurrente deba esperar a que la transacción actual confirme o aborte.
    
    \item \textbf{Prevención de Deadlocks:}
    Para evitar el problema de bloqueo mutuo (\textit{deadlocks}) donde dos transacciones concurrentes retienen recursos cruzados, se incorporó un algoritmo de ordenamiento determinista en la función de checkout. Antes de adquirir los bloqueos en base de datos, el backend ordena el array de productos por su UUID de forma ascendente. Al bloquear los recursos siempre en el mismo orden físico en el disco, se imposibilita la formación de ciclos de espera.
    
    \item \textbf{Daemon de Expiración (TTL) en Background:}
    En \texttt{backend/src/index.ts} se configuró un proceso daemon en segundo plano que corre cada 10 segundos. Este daemon escanea la base de datos de manera atómica buscando registros en la tabla \texttt{Reserva} que estén en estado \texttt{Pendiente} y cuya \texttt{fechaExpiracion} sea menor al tiempo actual. Al detectar expiraciones, abre una transacción que revierte las unidades reservadas devolviéndolas al stock vendible disponible y actualiza el estado a \texttt{Expirada}.
\end{enumerate}

\subsection{Evidencias de Funcionamiento y Logs}

A continuación se adjuntan las trazas de ejecución en consola obtenidas directamente del servidor de backend de Sync-Stock, demostrando el ciclo de vida de los servicios, el control de concurrencia y las tareas en segundo plano.

\paragraph{Logs de Inicialización del ESB y Conexión de Servicios:}
\begin{verbatim}
[ESB-TCP] Server listening on port 5000...
[ESB-TCP] New client connection established.
[ESB-TCP] Registered service: 'autho' on local socket.
[ESB-TCP] Registered service: 'inven' on local socket.
[ESB-TCP] Registered service: 'sales' on local socket.
[ESB-Mediator] ESB Mediator registering services via TCP sockets...
[ESB-Mediator] Services successfully mapped and listening for requests.
[Express-Gateway] Gateway API listening on port 3001.
\end{verbatim}

\paragraph{Logs de Transacción de POS y Bloqueo Pesimista (Checkout Concurrente):}
\begin{verbatim}
[Express-Gateway] POST /api/v1/sales/checkout received from POS
[ESB-TCP] Routing message to service 'sales' (action: 'checkout')
[Sales-Service] Processing checkout. Total items: 2.
[Inventory-Service] Initiating stock check and lock transaction.
[Inventory-Service] Sorted product list to prevent deadlocks:
  - SKU: PES-CANA-CARBON-01 (UUID: a1b2c3d4-...)
  - SKU: PES-SE-NZ-023 (UUID: f5b88220-...)
[Prisma-SQL] SELECT * FROM "Inventario" WHERE ... FOR UPDATE
[Prisma-SQL] Row locks acquired successfully. No contention.
[Inventory-Service] Stock verified. Deducting 2 units.
[Prisma-SQL] UPDATE "Inventario" SET "stockDisponible" = 13 WHERE ...
[Sales-Service] Order d3b07384-d113... recorded successfully.
[Inventory-Service] Transaction committed. Locks released.
[ESB-TCP] Transaction checkout response returned to Protocol Bridge.
[Express-Gateway] 201 Created returned to POS vendor.
\end{verbatim}

\paragraph{Logs del Daemon de Expiración Temporal de Reservas Web (TTL):}
\begin{verbatim}
[Daemon-TTL] Scanning database for expired web reservations...
[Daemon-TTL] Found 1 expired pending reservation (ID: b1b88220-432a...).
[Inventory-Service] Initiating reservation rollback transaction.
[Prisma-SQL] SELECT * FROM "Inventario" WHERE ... FOR UPDATE
[Inventory-Service] Reverting 2 units of SKU PES-CANA-CARBON-01.
[Prisma-SQL] UPDATE "Inventario" SET 
  "stockDisponible" = "stockDisponible" + 2,
  "stockReservado" = "stockReservado" - 2 WHERE ...
[Prisma-SQL] UPDATE "Reserva" SET "estado" = 'Expirada' WHERE ...
[Inventory-Service] Rollback transaction committed successfully.
[Daemon-TTL] Expired reservation b1b88220 reclaimed. Stock restored.
\end{verbatim}

\newpage
\section{Análisis Crítico}

\subsection{Latencia y Concurrencia de los Bloqueos Pesimistas}
El uso de bloqueos pesimistas (\texttt{SELECT ... FOR UPDATE}) garantiza una consistencia absoluta (consistencia estricta en el contexto del teorema de CAP), lo cual es crítico para erradicar las sobreventas en PescaTotal. Sin embargo, esta decisión de diseño conlleva un costo en latencia bajo escenarios de alta concurrencia. Cuando se bloquean filas de inventario para un producto altamente demandado (por ejemplo, ofertas flash en el canal web), las peticiones de checkout concurrentes se encolan secuencialmente en el motor PostgreSQL, aumentando la latencia general del sistema de $<50$ ms hasta picos de $>1000$ ms. Para mitigar esto en un entorno productivo de mayor escala, se debería evaluar una estrategia híbrida utilizando optimismo de stock mediante caches en memoria (ej. Redis) o reservas distribuidas desacopladas del motor de base de datos relacional.

\subsection{Desafíos del Desarrollo de un ESB Personalizado TCP}
La decisión de implementar un bus de servicios personalizado sobre sockets TCP crudos en Node.js permitió lograr un sistema ultraligero y sin dependencias de infraestructura externas. No obstante, en un entorno de desarrollo real esto introduce complejidades de bajo nivel que los brokers industriales (como RabbitMQ, Apache Kafka o gRPC) resuelven de forma nativa:
\begin{itemize}
    \item \textbf{Entramado de Datos y Fragmentación:} Al usar streams TCP crudos, los mensajes pueden llegar fragmentados o concatenados si se transmiten de forma muy veloz. Se tuvo que programar un algoritmo de parseo manual con prefijo de longitud para solucionar esto de forma segura.
    \item \textbf{Gestión de Colas y Reintentos:} A diferencia de RabbitMQ, nuestro bus customizado no almacena físicamente los mensajes en disco (encolamiento persistente) ni cuenta con mecanismos de Dead Letter Queues (DLQ) para reenvío de mensajes fallidos de forma automática.
\end{itemize}

\subsection{Seguridad y RBAC en Entornos Distribuidos}
El diseño e implementación de un middleware de control de acceso por roles (RBAC) centralizado en el Gateway REST protege con éxito la capa lógica del ESB de peticiones externas no autorizadas. Los usuarios del sistema solo interactúan con pantallas permitidas (por ejemplo, restringiendo la pestaña de Catálogo CRUD y la Auditoría de mensajería del ESB al administrador). Sin embargo, a nivel de red interna, la comunicación TCP entre el Protocol Bridge (Express) y el bus de servicios viaja actualmente en texto plano. En un entorno de producción distribuido real en la nube, es imperativo implementar encriptación de extremo a extremo mediante el módulo \texttt{tls} de Node.js, configurando certificados SSL/TLS autofirmados para mitigar ataques de escucha de red interna (Man-in-the-Middle).

\subsection{Control de Inactividad y Seguridad Física en Tienda}
La incorporación de la autodetección de inactividad de 30 minutos es un paso crítico para proteger la seguridad de la información física en la tienda PescaTotal. Dado que las terminales del Punto de Venta (POS) suelen compartirse entre diferentes vendedores en el mesón de atención física, un olvido de cierre de sesión podría permitir a un vendedor registrar ventas a nombre de otro o a un cliente acceder a configuraciones de stock del panel administrador. La simulación de inactividad implementada proporciona un mecanismo de auditoría rápido que valida que el estado reactivo del frontend se limpie correctamente al expirar la sesión.

\section{Conclusiones}
Con la definición del modelo de datos relacional y la estructuración en base a servicios (SOA), el sistema \textbf{Sync-Stock} adquiere la solidez técnica necesaria para erradicar la sobreventa en PescaTotal. La separación entre componentes cliente y servicio asegura que la plataforma sea escalable, permitiendo a futuro integrar nuevos canales de venta sin alterar el núcleo del control de inventario. La estricta vinculación de cada interfaz con los requerimientos de seguridad, consistencia y baja latencia asegura el cumplimiento tanto de los objetivos comerciales como de los estándares de ingeniería requeridos.
\end{document}